% Clase 23 --- El transformador (§3.1). Reemplaza el esquema inline de la
% primera version, restilado al vocabulario de tikzstyles.tex.
% Lo esencial que tiene que mostrar el dibujo: el nucleo cerrado obliga al MISMO
% flujo por espira en las dos bobinas; lo unico distinto es cuantas espiras ve.
\begin{center}
\begin{tikzpicture}[scale=0.85, transform shape]

  % ---- núcleo de hierro (marco cerrado) ----
  \fill[figgray!18] (0,0) rectangle (5,3.4);
  \fill[white] (1,0.7) rectangle (4,2.7);
  \draw[figgray, line width=0.9pt] (0,0) rectangle (5,3.4);
  \draw[figgray, line width=0.9pt] (1,0.7) rectangle (4,2.7);
  \node[figlblaux, anchor=center, align=center] at (2.5,1.7)
    {núcleo de\\ hierro};

  % ---- flujo confinado en el núcleo ----
  \draw[figamber, line width=1.1pt, -{Latex[length=2.4mm,width=1.8mm]}]
    (1.7,3.05) -- (2.7,3.05);
  \draw[figamber, line width=1.1pt, -{Latex[length=2.4mm,width=1.8mm]}]
    (2.7,0.35) -- (1.7,0.35);
  \node[figlbl, figamber, anchor=west] at (2.85,3.05) {$\Phi(t)$};

  % ---- primario ----
  \draw[figline] (0.5,0.9) to[cute inductor] (0.5,2.5);
  \draw[figline] (0.5,2.5) -- (0.5,4.0) -- (-2.2,4.0);
  \draw[figline] (0.5,0.9) -- (0.5,-0.6) -- (-2.2,-0.6);
  \draw[figline] (-2.2,-0.6) to[sV=$\mathcal{E}(t)$] (-2.2,4.0);
  \node[figlbl, anchor=east] at (0.15,1.7) {$N_p$};
  \node[figlbl, anchor=south, align=center] at (-0.9,4.1) {primario};

  % ---- secundario ----
  \draw[figline] (4.5,0.9) to[cute inductor] (4.5,2.5);
  \draw[figline] (4.5,2.5) -- (4.5,4.0) -- (7.2,4.0);
  \draw[figline] (4.5,0.9) -- (4.5,-0.6) -- (7.2,-0.6);
  \draw[figline] (7.2,4.0) to[R=$R_{\text{carga}}$] (7.2,-0.6);
  \node[figlbl, anchor=west] at (4.85,1.7) {$N_s$};
  \node[figlbl, anchor=south, align=center] at (5.9,4.1) {secundario};

\end{tikzpicture}\\[0.5em]
{\small\itshape El transformador no tiene ninguna conexión eléctrica entre los
dos circuitos: lo único que los vincula es el flujo magnético, y el núcleo de
hierro está ahí para que ese vínculo sea casi perfecto. Como
$\mu_{\text{hierro}}/\mu_{\text{aire}} \sim 4000$, el campo prefiere quedarse
dentro del marco cerrado en vez de escaparse al aire, y entonces
\textbf{cada espira} ---la primera del primario o la última del secundario---
ve el \emph{mismo} $\Phi(t)$. Ahí está el argumento completo: por Faraday cada
espira aporta $d\Phi/dt$, así que $V_p = N_p\,d\Phi/dt$ y
$V_s = N_s\,d\Phi/dt$, y al dividir se va lo único que no sabemos calcular,
quedando $V_p/V_s = N_p/N_s$. Nótese que todo esto exige corriente
\textbf{alterna}: con corriente continua $\Phi$ es constante, $d\Phi/dt = 0$ y
el secundario no ve absolutamente nada.}
\end{center}
